\documentclass[../../../../main.tex]{subfiles}
\begin{document}

\begin{flushright}
\footnotesize\textit{【自動文字起こし・要確認】}
\end{flushright}

\begin{proof}[解]
$U_1 \dots 4n$ 個の赤 と $n$ 個の白\\
$U_2 \dots 2n$ 個の赤 と $3n$ 個の白

誤った判断を下すのは、「$U_1$ が手渡されて $U_2$ と判断する $\dots A$」か「$U_2$ が手渡されて $U_1$ と判断する $\dots B$」のいずれか.

\medskip
1° Aの時\\
$U_1$ から $(\text{赤}, \text{白}) = (1, 2), (0, 3)$ をとり出した時で 確率は
\[
\frac{2}{3} \times \frac{_{4n}\mathrm{C}_1 \cdot _n\mathrm{C}_2 + _n\mathrm{C}_3}{_{5n}\mathrm{C}_3}
\]
(すべての玉を区別し、そのうちから3つをとり出す $_{5n}\mathrm{C}_3$ 通りが同様に確からしい)

\medskip
2° Bの時\\
$U_2$ から $(\text{赤}, \text{白}) = (3, 0), (2, 1)$ をとり出した時で 確率は
\[
\frac{1}{3} \times \frac{_{2n}\mathrm{C}_3 + _{2n}\mathrm{C}_2 \cdot _{3n}\mathrm{C}_1}{_{5n}\mathrm{C}_3}
\]

\medskip
以上から
\[
P_n = \frac{1}{3 \cdot _{5n}\mathrm{C}_3} \left[ 2 \cdot _{4n}\mathrm{C}_1 \cdot _n\mathrm{C}_2 + 2 \cdot _n\mathrm{C}_3 + _{2n}\mathrm{C}_3 + _{2n}\mathrm{C}_2 \cdot _{3n}\mathrm{C}_1 \right]
\]
\[
= \frac{6}{15n(5n-1)(5n-2)} \cdot \frac{1}{6} \left[ 24 n^2(n-1) + n(n-1)(n-2) + 2n(2n-1)(2n-2) + 18n^2(2n-1) \right]
\]
\[
= \frac{1}{15} \cdot \frac{1}{(5 - \frac{1}{n})(5 - \frac{2}{n})} \left[ 24\left(1 - \frac{1}{n}\right) + \left(1 - \frac{1}{n}\right)\left(1 - \frac{2}{n}\right) + 2\left(2 - \frac{1}{n}\right)\left(2 - \frac{2}{n}\right) + 18\left(2 - \frac{1}{n}\right) \right]
\]
\[
\longrightarrow \frac{1}{15} \cdot \frac{1}{25} [ 24 + 1 + 8 + 36 ] = \frac{69}{375} \quad (n \to \infty)
\]
\end{proof}

\end{document}
